\section{De dónde vienen los datos: Web, Code, Curated y Synthetic}

Una vez fijado el token budget, la segunda pregunta es la source. Las distintas fuentes no solo difieren en contenido: también implican distintos costos de rastreo, licencia, deduplicación, formato y verification. La clase divide las fuentes en web data, GitHub/code, curated sources y synthetic data; StarCoder, por su parte, muestra cómo convertir un code corpus en un producto de datos que puede inspeccionarse.

\subsection{Source taxonomy: la fuente determina los costos de gobernanza posteriores}

La sección anterior solo estimó el token budget; esta sección empieza a responder de dónde vienen esos tokens, porque la source determina los costos posteriores de filter, license y evaluation. La web ofrece amplia cobertura, pero con mucho ruido; GitHub tiene estructura y metadatos de licencia, pero contiene duplicados y secrets; las curated sources son de alta calidad, pero de escala limitada; la synthetic data puede generarse de forma dirigida, pero depende del teacher, el prompt y la seed diversity.

\lecturefigure{slide-26.jpg}{Las cuatro categorías de fuentes de LLM training data}{Deck oficial de Loubna Ben Allal, página 26}

\begin{knowledgebox}{Las fuentes no son etiquetas mutuamente excluyentes}
Los repositorios de código contienen README/Markdown, las páginas web en lenguaje natural contienen código, y los synthetic prompts a menudo se derivan de web/curated documents. La mixture real debe registrarse según la provenance y la transformation lineage, no solo con los porcentajes finales por domain.
\end{knowledgebox}

\subsection{Common Crawl: una escala tan grande que obliga a filtrar antes de descargar/procesar}

Common Crawl ofrece instantáneas periódicas de la web; el raw corpus puede alcanzar cientos de TB. Tokenizarlo directamente no solo es costoso, sino que conserva boilerplate, spam, machine-generated pages, duplicates, PII y texto de baja calidad; por eso, normalmente se reutiliza un filtered web dataset existente o se aplica distributed filtering.

\lecturefigure{slide-27.jpg}{Web data: la escala de Common Crawl y la carga de filtrado}{Deck oficial de Loubna Ben Allal, página 27}

\begin{warningbox}{Rastrear una página web no equivale a tener permiso para entrenar con ella}
La accesibilidad, robots, copyright, terms, personal data y jurisdiction son cuestiones distintas. Un pipeline técnico puede registrar la procedencia y retirar muestras, pero no puede resolver automáticamente todos los juicios legales y éticos.
\end{warningbox}

\teachervoice{00:13:20--00:16:20, la ponente subraya la escala original de Common Crawl y el filtering cost, y recomienda entender el procesamiento de los datasets existentes en lugar de tratar el raw crawl como datos listos para usarse.}

\subsection{FineWeb: un mejor filter se demuestra con la early-training curve}

La slide de FineWeb compara las curvas de entrenamiento de modelos pequeños sobre varios web corpora. Al leer la figura, primero hay que confirmar que la architecture, el token budget, la evaluation y la seed sean los mismos; después, observar qué corpus alcanza antes un score más alto. Lo que la curva respalda es que el filter pipeline supera al control en ese entorno, no que sea el mejor para cualquier modelo/tarea.

\lecturefigure{slide-29.jpg}{FineWeb ablation: early-training performance de distintos web corpora}{Deck oficial de Loubna Ben Allal, página 29}

\begin{importantbox}{La operational definition de data quality}
Alta calidad no significa «parecer un buen artículo». En el preentrenamiento, debe manifestarse como menor validation loss, downstream learning más rápido, menos memorization/toxicity/PII y mejor domain coverage, y debe sostenerse con matched compute y varias seeds.
\end{importantbox}

\subsection{The Stack: los datos de código requieren raw collection, inspection y opt-out}

El código no es como una página web de texto plano: los archivos tienen estructura (repository, path, license, language, stars, commit history, etc.) y también pueden incluir generated files, vendor copies, secrets y benchmark solutions. Por eso, el pipeline de The Stack no solo recolecta: también realiza detección de lenguaje, gestión de license/metadata, dedup e inspection por parte de la comunidad.

\lecturefigure{slide-30.jpg}{Code data: panorama de collection y processing de The Stack}{Deck oficial de Loubna Ben Allal, página 30}

\lecturefigure{slide-31.jpg}{The Stack v1: GitHub collection, raw data, dedup, inspection y opt-out}{Deck oficial de Loubna Ben Allal, página 31}

\begin{knowledgebox}{El lugar del opt-out en el pipeline}
El opt-out requiere stable repository/file identifiers, herramientas públicas de consulta, una lista versionada de las solicitudes procesadas y un mecanismo de exclusión para el future retraining. No consiste en borrar un enlace web una vez terminado el entrenamiento, sino que forma parte de la data lineage.
\end{knowledgebox}

\teachervoice{00:17:10--00:19:10, la ponente explica que BigCode publicó con la v1 una inspection tool y un opt-out form, para que los desarrolladores pudieran revisar el código y solicitar su exclusión de entrenamientos futuros; este es un diseño central de los proyectos de datos abiertos.}

\subsection{The Stack v2: Software Heritage y más datos del proceso de desarrollo}

La v2 ya no depende principalmente de clonar GitHub de forma directa: extrae los datos del Software Heritage archive e incorpora issues, pull requests, commits, Jupyter/Kaggle notebooks y math/coding datasets. Así, el corpus no solo enseña el código final, sino que también contiene los problemas, las modificaciones, las discusiones y los entornos de ejecución.

\lecturefigure{slide-32.jpg}{The Stack v2: Software Heritage, filtrado y datos de desarrollo de múltiples fuentes}{Deck oficial de Loubna Ben Allal, página 32}

\lecturefigure{slide-33.jpg}{The Stack v1 y v2: raw size, filtered size, languages y tokens}{Deck oficial de Loubna Ben Allal, página 33}

\begin{importantbox}{La brecha entre raw scale y usable scale es trabajo de ingeniería}
El raw size de la v2 es aproximadamente diez veces el de la v1, y después del filtrado sigue siendo unas cuatro a cinco veces el de la v1. La gran mayoría de lo eliminado no es «desperdicio», sino duplicates, contenido de baja calidad, unsupported formats, PII o lenguajes no adecuados; la tasa de filtrado por sí sola no indica si el resultado es bueno o malo: hay que observar la diversity conservada y los resultados del modelo.
\end{importantbox}

\teachervoice{00:19:10--00:20:20, la ponente sitúa las mejoras centrales de la v2 en la archive source, más lenguajes y datos auxiliares de desarrollo, y no solo en «más tokens».}

\subsection{Synthetic data: de Phi a un corpus de generación controlable}

La serie Textbooks Are All You Need/Phi muestra que, con datos de alta calidad con estilo de libro de texto generados por GPT-3.5/4, es posible entrenar modelos pequeños potentes con un corpus mucho menor que uno de web-scale. Lo que esto suscitó no fue la conclusión de que «lo synthetic sustituye a los datos reales», sino que generator, prompt, seed y verification se convirtieran en nuevas variables de la ingeniería de datos.

\lecturefigure{slide-34.jpg}{Caso catalizador de la synthetic data: Textbooks Are All You Need}{Deck oficial de Loubna Ben Allal, página 34}

Sea $s$ la seed/context, $p$ el prompt template y $G$ el generator; la synthetic sample es
\[
y\sim G(\cdot\mid p,s).
\]
La calidad de los datos depende de la cobertura de $s$, las restricciones de $p$, la capacidad de $G$ y el verifier a posteriori $v(y)$; no puede depender solo de la cantidad de muestras.

\begin{warningbox}{Synthetic feedback loop}
Si el generator y el verifier comparten puntos ciegos, la generación repetida amplifica la homogeneización del estilo, los errores factuales y el representation collapse. Conviene registrar teacher/version, seed provenance, rejection reasons, novelty y human spot checks.
\end{warningbox}

\subsection{Cosmopedia: ampliar la diversidad con web/curated seeds}

Cosmopedia usa Mixtral-8x7B para generar unos 25B tokens de libros de texto, blogs, cuentos, etc. La técnica central no es una simple instrucción del tipo «haz que el modelo escriba un libro de texto», sino extraer context de web samples, Stanford courses, WikiHow, etc., y colocar en el prompt el topic junto con una seed concreta, para limitar la hallucination y aumentar la dispersión temática.

\lecturefigure{slide-35.jpg}{Cosmopedia: synthetic text dataset a gran escala generado con modelos abiertos}{Deck oficial de Loubna Ben Allal, página 35}

\lecturefigure{slide-36.jpg}{Composición de los datos de Cosmopedia: web seeds, curated sources y múltiples géneros}{Deck oficial de Loubna Ben Allal, página 36}

\begin{knowledgebox}{Las cuatro auditorías de un synthetic corpus}
\begin{enumerate}
  \item \textbf{Coverage:}¿están dispersos el topic, el language, la difficulty y el style?
  \item \textbf{Grounding:}¿la generación está restringida por la seed/context?
  \item \textbf{Verification:}¿se revisan de forma automática/humana los hechos, el código y las matemáticas?
  \item \textbf{Lineage:}¿pueden rastrearse teacher, prompt, seed, filter y license?
\end{enumerate}
\end{knowledgebox}

\teachervoice{00:20:20--00:24:20, la ponente dice que la synthetic data es muy reciente y muy prometedora, pero que obtener diversity es «very tricky». Cosmopedia usa una gran cantidad de web/curated seeds, en lugar de ampliar sin límite a partir de unos pocos prompts vagos.}

\subsection{Resumen de la sección}

La source de los datos determina los costos de procesamiento posteriores. La web aporta cobertura; el code, estructura; las curated sources, una signal alta; y la synthetic data, ampliación dirigida. Cualquier fuente debe ir acompañada de un diseño de provenance, filter, license/opt-out y evaluation.
